% !PRINTED 28
\weberpage{28}
%<<CONT
\begingroup
\setlength{\abovedisplayskip}{5pt plus 2pt minus 1pt}
\setlength{\belowdisplayskip}{5pt plus 2pt minus 1pt}
\setlength{\abovedisplayshortskip}{3pt plus 1pt}
\setlength{\belowdisplayshortskip}{4pt plus 1pt}
\linespread{0.94}\selectfont
that $m\geqq n$ and that $a_0$ and $b_0$ are different from zero. Then
the difference
\begin{equation}
 A-{a_0\over b_0}x^{m-n}B
\tag{2}
\end{equation}
is also a whole rational function of $x$, but its degree is
\gesperrt{smaller} than $m$, since the highest power in the two terms
of the difference has the same coefficient and therefore cancels. We
put this difference equal to
\begin{equation}
 A'=A'(x)=a'_0x^{m'}+a'_1x^{m'-1}+\cdots,
 \qquad m'<m.
\tag{3}
\end{equation}
If now $m'$ is still $\geqq n$, then in (2) we can put $A'$ in place of
$A$ and repeat the same reasoning.

This gives a chain of equations:
\begin{equation}
\begin{aligned}
 A-{a_0\over b_0}x^{m-n}B&=A',\\
 A'-{a'_0\over b_0}x^{m'-n}B&=A'',\\
 A''-{a''_0\over b_0}x^{m''-n}B&=A''',\\
 &\cdots,
\end{aligned}
\tag{4}
\end{equation}
and this chain can be continued until the degree of the function that
has arisen is smaller than $n$. Since in the sequence of functions
$A,A',A''\ldots$ the degree of each following function is lowered by
\gesperrt{at least} one unit, the chain of equations (4) consists of
\gesperrt{at most} $m-n+1$ members; but it may also contain fewer
members when \gesperrt{several} powers cancel simultaneously. If we add
all equations (4), denote the last of the functions $A,A'\ldots$ by
$C$, and, for abbreviation, set
\begin{equation}
 Q={a_0\over b_0}x^{m-n}+{a'_0\over b_0}x^{m'-n}+\cdots,
\tag{5}
\end{equation}
so that $Q$ also is a whole rational function of $x$, then there
follows
\begin{equation}
 A=QB+C.
\tag{6}
\end{equation}
The operation just described, by which the functions $Q,C$ are found
from $A,B$, is called \gesperrt{division}. $A$ is the
\gesperrt{dividend}, $B$ the \gesperrt{divisor}, $C$ the
\gesperrt{remainder}, and $Q$ the \gesperrt{quotient}. \gesperrt{The
degree of the remainder is always lower than the degree of the
divisor.}
\endgroup
